\section{在行动时刻之前完成计划}

\begin{proposition}[已解决的进入会降低含糊性]
如果一个计划为目标状态 $s_1$ 提供了已解决的进入规格，那么与一个在其他方面相似、但没有计划的意图相比，目标含糊项 $U(s_1)$ 会弱降低。
\end{proposition}

\begin{proof}
根据定义，已解决的进入规格固定了一个可见的第一步动作，而若没有它，这一步动作本会在行动时刻仍然悬而未决。由于 $U(s_1)$ 被定义为进入 $s_1$ 时尚未解决的规格负担，解决至少一个这样的负担，就会使 $U(s_1)$ 弱降低。
\end{proof}

\begin{proposition}[在单调支架下，计划会降低未来障碍]
假设一个关于 $s_0 \to s_1$ 的计划同时降低了目标含糊性并提高了相关准备。在单调计划假设下，未来的激活障碍将弱低于没有该计划的情形。
\end{proposition}

\begin{proof}
由前一个命题可知，已解决的进入会降低 $U(s_1)$。按照计划的定义，相关准备也会增加。单调计划假设说明，$\DV(s_0 \to s_1)$ 关于含糊性弱增，关于准备弱减。因此，降低含糊性并增加准备，会使激活障碍弱降低。
\end{proof}

\begin{corollary}[临时自我说服在结构上更弱]
如果临时自我说服既不改变含糊性，也不改变准备，而计划至少会在一个方向上有利地改变其中一项，那么带计划的转变，其激活障碍就会弱低于仅被说服的转变。
\end{corollary}

\begin{proof}
根据假设，仅被说服的转变不会改变相关障碍变量。带计划的转变则至少降低一个有利项。应用前一个命题即可。
\end{proof}

\begin{remark}
这个推论刻意比本章早先版本中的“计划定理”更弱。它并不主张计划在所有人类情境中都一定获胜。它只是在说：如果计划真的改变了模型所认为重要的那些变量，那么模型就会预测，计划比单纯的口头认可拥有更低的障碍。
\end{remark}